\section{Objetivo de costos}

El objetivo de costos busca medir la eficiencia en el uso de los recursos asociados a la operación de la flota, permitiendo identificar oportunidades de optimización y reducción de gastos operativos. A través de este indicativo se evalúa el desempeño de los vehículos en términos de consumo y utilización.

\subsection{Rendimiento de combustible}

Este indicador mide la eficiencia del consumo de combustible de los vehículos, relacionado a la distancia recorrida con la cantidad de combustible consumido durante el periodo analizado. Su principal objetivo es evaluar el desempeño operativo de la flota y detectar posibles oportunidades de mejora en los hábitos de conducción.

\begin{equation*}
    \text{rendimiento combustible} = \sum_i^N \frac{\text{km recorridos}}{\text{consumo de combustible}} \cdot W\big( \hspace{0.07cm} \smash{\vec{t}_i}\,\vphantom{t} \hspace{0.07cm} \big)
\end{equation*}

Un valor alto indica una mayor eficiencia en el consumo de combustible, ya que el vehículo recorre una mayor distancia por cada galón consumido.
